\section{策略梯度的实现与调试}
\subsection{伪代码与最佳实践}
下面是典型的 policy gradient 伪代码：
\begin{enumerate}
    \item 启动环境，初始化 $\theta$，准备 buffer。
    \item 收集 $N$ 条 trajectory，记录 $(s_t,a_t,r_t)$。
    \item 计算 reward-to-go 与 baseline（可用 value network）。
    \item 构造 surrogate loss，执行一轮 gradient update。
    \item 重复，每 $K$ 步 eval 一次策略。
\end{enumerate}

\begin{table}[H]
\centering
\small
\begin{tabular}{p{0.32\textwidth}p{0.4\textwidth}p{0.26\textwidth}}
\toprule
组件 & 建议范围 & 说明 \\
\midrule
Batch size & 2048--8192 steps & 决定 variance 和 compute \\
Learning rate & 1e-4 -- 5e-4 & 小 learning rate 保证稳定 \\
Entropy coeff & 0.01--0.03 & 激励探索，防止收敛太快 \\
$\lambda$ & 0.9--0.98 & GAE smoothing \\
\bottomrule
\end{tabular}
\caption{训练策略梯度时的关键超参数}
\end{table}

\subsection{调试提醒}
\begin{warningbox}{忽略日志会错过差错}
策略梯度容易发生梯度爆炸/消失、exploding reward scale、policy collapse。必须记录 reward curve、policy entropy、KL divergence，否则难以定位问题。
\end{warningbox}

\begin{importantbox}{调试建议}
\begin{itemize}
    \item 先从小环境（CartPole）验证数据 pipeline；
    \item 手动监控 advantage 分布，确保其均值在零附近；
    \item 使用 gradient clipping 或 trust region 避免大步更新。
\end{itemize}
\end{importantbox}

\subsection{本章小结}
实现层面需要把采样、估计、优化、日志串成闭环；合理的超参、调试指标与记录机制是防止训练崩坏的保障。

